% Supplementary practice source only; outside the 100-mark core paper.
% Requires the shared midterm style when included in another document.
\newpage
\textbf{\Large แบบฝึกเสริมนอกคะแนน}\par
S1-S3 ใช้ฝึกประเภทที่ไม่ได้เลือกในชุดหลัก ไม่รวมใน 100 คะแนนหรือเวลา 3 ชั่วโมง ข้อ 1-12 ยังคงเป็นชุดจำลองสอบเดิม\par
\medskip
\textbf{S1\quad Use Case Diagram: University Grades Submission}\par
\textbf{ข้อมูลสำหรับทำโจทย์} (เรียบเรียงจาก Lecture 4 หน้า 64-65)\par
\begin{enumerate}[itemsep=3pt]
\item Instructor เข้าระบบด้วย KU account ผ่าน KU Single Sign-On เฉพาะบัญชีที่ถูกต้องเท่านั้นที่เข้าใช้งานได้
\item Instructor เลือก semester แล้วดู courses และ sections ที่เปิดใน department ของตนจาก KU Registration System
\item เมื่อเลือก course และ section(s) ระบบดึงรายชื่อนักศึกษาที่ลงทะเบียนจาก KU Registration System มาใส่ grade sheet ได้
\item Grade sheet มี header (Course, Semester) และรายชื่อนักศึกษา เริ่มใน Draft Instructor กรอกหรือแก้เกรด และเลือกพิมพ์สำเนาเก็บไว้ได้
\item เมื่อกรอกเกรดครบ Instructor เปลี่ยน sheet เป็น Ready แล้วจึง submit ได้; ห้าม submit จาก Draft
\item เมื่อ submit ระบบอัปเดต transcript ของนักศึกษาทุกคน แจ้งเกรดแต่ละคนทาง email ยืนยันผลสำเร็จ และเปลี่ยน sheet เป็น Submitted
\end{enumerate}
\textbf{ระบบภายนอกทั้งห้า}\par
\begin{tabularx}{\linewidth}{|>{\raggedright\arraybackslash}m{5.3cm}|Y|}\hline
KU Single Sign-On & ตรวจบัญชีและจัดการรหัสผ่าน\\\hline
KU Registration System & ให้ข้อมูล courses/sections และผู้ลงทะเบียน\\\hline
KU Transcript System & บันทึกเกรดอย่างเป็นทางการ\\\hline
KU Student Information System & ให้ข้อมูลติดต่อ เช่น email ของนักศึกษา\\\hline
KU Email Service (Google Workspace) & ส่ง email แจ้งเกรด\\\hline
\end{tabularx}
\par\medskip
\textbf{คำสั่ง:} วาด Use Case Diagram ในหน้าถัดไปให้ครอบคลุมเป้าหมายของ Instructor และระบบภายนอก แสดง system boundary, actors และ associations ให้ถูก หากข้อกำหนดรองรับ ให้แสดง \code{<<include>>}, \code{<<extend>>} หรือ generalization ด้วยทิศทางและสัญลักษณ์ที่ถูกต้อง; ระบุเงื่อนไขของ \code{<<extend>>} และอย่าสร้างความสัมพันธ์ที่โจทย์ไม่ได้กำหนด\par
\textbf{อธิบายสั้น ๆ:} อธิบายความสัมพันธ์ที่ใช้ว่าเป็นขั้นบังคับหรือทางเลือก เงื่อนไขใดทำให้ submit ได้ และมี generalization ที่ข้อกำหนดรองรับหรือไม่\par
\blankbox{3.5cm}
\newpage
\textbf{\Large S1 (ต่อ): พื้นที่วาด Use Case Diagram}\par
University Grades Submission\hfill แบบฝึกเสริม ไม่รวม 100 คะแนน\par
\blankbox{21.5cm}
\newpage
\textbf{\Large S2\quad State Diagram: Service Ticket}\par
แบบฝึกเสริม ไม่รวม 100 คะแนน ใช้เส้นทางทำงานของตัวอย่าง Lecture 4 หน้า 61-62 และข้อมูลจากหน้า 45-46\par
\begin{enumerate}[itemsep=2pt]
\item เมื่อบันทึก Ticket ใหม่ System สร้าง Action Record แรกโดยอัตโนมัติ กำหนด status เป็น Requested และ New Due Date เท่ากับ Initial Due Date
\item CCE ทำ action ได้เฉพาะ Requested: อนุมัติเป็น Approved หรือปิดเป็น Closed by CCE พร้อม notes/วันนัดใหม่ได้ตามกรณี
\item ตามเส้นทางตัวอย่าง Programmer เริ่มงานจาก Approved ให้เป็น In Progress แล้วเพิ่ม progress actions เพื่อบันทึก notes หรือวันนัดใหม่ได้โดยยังเป็น In Progress
\item เมื่อทำงานเสร็จ Programmer ปิดเป็น Closed by Programmer และกำหนด New Due Date เป็น completion date
\item หลังบันทึกทุก action ระบบอัปเดต Last Status และ Last Due Date ของ Ticket จาก action ล่าสุด; action ที่บันทึกแล้วแก้ไขหรือลบไม่ได้
\item Ticket ที่ Closed by CCE หรือ Closed by Programmer เพิ่ม action ไม่ได้อีก Customer และ IT Admin ไม่ได้ทำ actions สำหรับเปลี่ยนสถานะเหล่านี้
\end{enumerate}
\textbf{คำสั่ง:} วาดเฉพาะวงจรสถานะของ Ticket หนึ่งใบ แสดง initial/final nodes, states, transitions และ self-loop ระบุชื่อ event และ actor ในวงเล็บเหลี่ยมตามแนว เช่น \code{approve [CCE]} ไม่ใช้ Customer/CCE/Programmer เป็นชื่อ state และไม่เติม state ที่โจทย์ไม่กำหนด\par
\blankbox{13cm}
\newpage
\textbf{\Large S3\quad Git: อธิบาย repository graph}\par
แบบฝึกเสริม ไม่รวม 100 คะแนน เรียบเรียงจาก Samples ข้อ 2.2 หน้า 8-9 ใช้ ExamScore และข้อสมมติ Git เดียวกับข้อ 11-12 ทุกคำถามอ่านจาก snapshot นี้แยกกัน เว้นแต่ระบุให้รันคำสั่ง\par
\begin{center}\supplementgitgraph\end{center}
\textbf{S3.1} local branch ใดถูก checkout และ HEAD ระบุ commit ใด?\par
\blankbox{1.5cm}
\textbf{S3.2} commit ใดบนปลายสาย local sprint ยังไม่อยู่ในประวัติที่ origin/sprint\_1\_staging อ้างถึง? อธิบายจากตัวชี้\par
\blankbox{1.8cm}
\textbf{S3.3} ไฟล์ใน working directory ตอนนี้มาจาก commit ใด?\par
\blankbox{1.5cm}
\textbf{S3.4} ถ้ารัน \code{git merge feature/load\_scores} ตอนนี้ fast-forward ได้หรือไม่? ต้องรวมแบบใดและเพราะอะไร?\par
\blankbox{2cm}
\textbf{S3.5} graph พิสูจน์ได้หรือไม่ว่า merge จะเกิด file conflict? เพราะอะไร?\par
\blankbox{1.5cm}
\textbf{S3.6} \code{git fetch origin} อาจอัปเดตอะไร และจะไม่เลื่อนอะไรให้โดยอัตโนมัติ?\par
\blankbox{2cm}
